% 1-3-questions.tex
%  Budget: 1pp.
%  Rules: every numeral in \lr{}; cite only from references.bib;
%  one unit at a time -- see thesis/WRITING_HANDOFF.md and APPROVALS.md.
%
%  AUTHOR CHECK REQUIRED BEFORE SUBMISSION: RQ1-RQ3 below are the working set
%  adopted in WRITING_HANDOFF.md §7 (2026-08-31), reconstructed from the work
%  actually done. They must be checked word-for-word against the approved
%  university proposal. RQ4 is not provisional. See CHECKLIST.md P0.

\section{سؤالات پژوهش}

پژوهش حاضر به چهار پرسش زیر پاسخ می‌دهد. سه پرسش نخست به مقایسه، ترکیب و
تفسیر مدل‌ها می‌پردازند و پرسش چهارم به انتخاب مسیر پیش‌بینی هدف روزانه
مربوط است.

\subsection*{پرسش نخست \lr{---} دقت مقایسه‌ای خانواده‌های مدل}

در پیش‌بینی روز-پیشِ قیمت برق بازار \lr{EPEX-DE}، پنج مدل منتخب \lr{---} ساده،
\lr{SARIMAX}، \lr{LEAR-LASSO}، \lr{LightGBM} و \lr{LSTM} \lr{---} تحت یک پروتکل
اعتبارسنجی پیش‌رونده یکسان، بر پایه معیارهای \lr{MAE}، \lr{RMSE}، \lr{sMAPE}
و \lr{rMAE}، چه عملکردی نسبت به یکدیگر دارند، و کدام تفاوت‌ها بر اساس آزمون
دیبولد-ماریانو با تصحیح \lr{HAC} از نظر آماری معنادار است؟

\subsection*{پرسش دوم \lr{---} ترکیب مدل‌ها و وزن‌دهی مشروط به رژیم}

آیا ترکیب وزن‌دار مدل‌های پایه دقت را نسبت به بهترین تک‌مدل بهبود می‌دهد، و
آیا مشروط‌کردن وزن‌ها به رژیم بازار \lr{---} که با آستانه‌ای صرفاً بر پایه داده
اعتبارسنجی به دو حالت آرام و پرتنش تفکیک می‌شود \lr{---} بهبود بیشتری نسبت به
وزن‌دهی ایستا ایجاد می‌کند؟ این بهبود، در صورت وجود، تحت کدام شرایط بازار
پدید می‌آید؟

\subsection*{پرسش سوم \lr{---} تفسیرپذیری}

بر پایه مقادیر \lr{SHAP} محاسبه‌شده برای مدل \lr{LightGBM}، کدام خانواده‌های
ویژگی بیشترین سهم را در پیش‌بینی دارند، و آیا این سهم میان روزهای آرام و
پرتنش، و میان افق ساعتی و روزانه، تغییر می‌کند؟

\subsection*{پرسش چهارم \lr{---} مسیر مستقیم در برابر مسیر تجمیعی}

برای پیش‌بینی میانگین بار پایه روزانه، کدام مسیر دقت بالاتری دارد: مدل‌سازی
مستقیمِ هدف روزانه، یا میانگین‌گیری از \lr{24} پیش‌بینی ساعتی؟ و آیا این
نتیجه برای هر پنج مدل یکسان است؟

\subsection*{یادداشتی درباره دامنه پرسش‌ها}

هر چهار پرسش عمداً به شواهدی محدود شده‌اند که این پژوهش در عمل تولید کرده
است. هیچ‌یک از آن‌ها ادعایی درباره برتری عمومی یک خانواده مدل در همه
بازارها مطرح نمی‌کند؛ دامنه اعتبار پاسخ‌ها یک بازار، یک بازه آزمون مشخص و یک
پروتکل ارزیابی واحد است. حدود دقیق این ادعاها در بخش \lr{4-7} و محدودیت‌های
آن‌ها در فصل پنجم بیان می‌شود.

آزمون برون‌توزیع مدل‌های تثبیت‌شده بر داده زنده سال \lr{2026}، که پس از
تصویب پیشنهاد پژوهش به کار افزوده شد، ذیل هیچ‌یک از این چهار پرسش قرار
نمی‌گیرد و در فصل پنجم به‌عنوان یافته‌ای مستقل و در عین حال محدودیتی جدی
گزارش می‌شود.
